% 原文 pp.178--187
% 原文 p.178
\hypertarget{srcpage:178}{}\label{srcpage:178}
\subsectionlink{量子Brown運動}
量子Brown運動の模型は、空間の一次元を運動する粒子が、温度$T$で熱平衡にある独立な調和振動子の環境と、線形に相互作用するものである。量子Brown運動は、おそらく最も重要で、最もよく研究されているデコヒーレンス模型である。特に日常世界の巨視的尺度でデコヒーレンス効果を見積もる際の標準的な選択となるだけでなく、散逸的な量子系の一般的研究でも、きわめて重要な役割を果たす\cite{ref018,ref218}。したがって、この量子Brown運動を、少し時間をかけて調べることにする。

以下の順に進む。まず5.2.1節では、第4章で展開した形式を用い、この模型のBorn--Markovマスター方程式を導く。5.2.2節では、中心粒子が調和振動子ポテンシャルに閉じ込められる特別な場合のマスター方程式を論じる。5.2.3節では、環境の性質を記述する、いわゆるスペクトル密度の特定の形に注目し、得られたマスター方程式をさらに評価する。5.2.4節では、高温環境という特別な場合を考え、デコヒーレンスの研究で中心的役割を果たす、対応するCaldeira--Leggettマスター方程式を導く。5.2.5節では、位相空間の二つのGaussian波束の重ね合わせの時間発展を調べ、量子Brown運動の動力学を解析する。環境との相互作用によって、位相空間で局在した優先状態が出現する仕組みを示す。これらの指針状態の動的な発展は、Newton軌道の「量子版」に対応する。Caldeira--Leggett模型（および量子Brown運動の一般的模型）の結果を過度に一般化することへの注意を、5.2.6節でいくつか述べる。最後の5.2.7節では、量子Brown運動の厳密なマスター方程式を簡単に論じる。

\subsubsectionlink{Born--Markovマスター方程式の導出}
まず模型を定義しよう。通常どおり、系と環境の全Hamiltonian $\hat H$を三つの部分に分解する［(2.85)、(4.12)参照］。
\chfiveeq{2}
当面、系の自己Hamiltonian $\hat H_{\mathrm S}$は指定しない（後の5.2.2節では調和振動子ポテンシャルの場合に限定する）。環境の自己Hamiltonian $\hat H_{\mathrm E}$は、調和振動子の集まりを記述する。

% 原文 p.179
\hypertarget{srcpage:179}{}\label{srcpage:179}
仮定により、これらの振動子は互いに相互作用しないので、$\hat H_{\mathrm E}$は単に個々の振動子のHamiltonianの和、すなわち
\chfiveeq{3}
となる。明らかな記法で、$m_i$、$\omega_i$は第$i$環境振動子の質量と固有振動数、$\hat q_i$、$\hat p_i$は正準位置演算子と運動量演算子を表す。

系と環境の相互作用の形は、中心粒子（系S）の位置座標$\hat X$が、環境の振動子の位置$\hat q_i$に、結合強度$c_i$で線形に結合するものである。したがって、相互作用Hamiltonian $\hat H_{\mathrm{int}}$は
\chfiveeq{4}
と書ける。$\hat X$、$\hat q_i$の両方に線形である双線形結合の仮定は、模型の重要な仮定であり、このため、しばしば「線形量子Brown運動」と明示的に呼ばれる。以下の議論は、最も一般的に研究される、この模型に基づく。

系と環境の相互作用は、環境が系の位置座標を連続的に監視することを記述する、(2.91)の形をしている。したがって、さらに計算をしなくても、結合(5.4)は、散逸効果や他の基底でのデコヒーレンスなど、ほかに生じうる帰結とともに、系の位置基底でのデコヒーレンスをもたらすと、すでに予想できる。

最初の目標は、環境自己相関関数(4.10)の決定である。$\hat H_{\mathrm{int}}$［(5.4)参照］には一項しかないので、計算すべき関数も一つだけ、すなわち\footnote{前の第4章と同様、本章全体を通じて、密度演算子以外の、陽な時間引数をもつ演算子は、すべて相互作用描像の演算子と理解する約束を用いる。}
\chfiveeq{5}
である。ここで$\hat\rho_{\mathrm E}=\hat\rho_{\mathrm E}(0)$［(4.6)参照］である。(5.4)を用いて評価すると
\chfiveeq{6}
を得る。表記を簡単にするため、ここでは、環境の異なる部分に属する演算子の積におけるテンソル積記号「$\otimes$」を省略した。本章を通じて、引き続き省略する。

% 原文 p.180
\hypertarget{srcpage:180}{}\label{srcpage:180}
(5.6)の最後の段階で$i\ne j$の項が消えるのは、環境振動子が互いに相互作用せず、したがって完全に無相関だからである。したがって$i\ne j$に対して
\chfiveeq{7}
となる。振動子の位置座標の期待値$\langle q_i\rangle_{\hat\rho_{\mathrm E}}$がゼロであることを用いた。

こうして、$C(\tau)$を評価する作業は、平均$\langle\hat q_i(\tau)\hat q_i\rangle_{\hat\rho_{\mathrm E}}$［(5.6)参照］を計算する作業に帰着した。これは容易に行える。$\hat q_i$を、ボソンの生成・消滅演算子$\hat a_i^{\dagger}$、$\hat a_i$で表す表示へ移ろう。
\chfiveeq{8}
すると、相互作用描像での$\hat q_i$の時間発展は、［(A.4)参照］
\chfiveeq{9}
と書ける。したがって
\chfiveeq{10}
である。しかし量
\chfiveeq{11}
は、単に環境の第$i$振動子の平均占有数にすぎない。仮定により環境は熱平衡にあり、これは
\chfiveeq{12}
に対応する。この表式と、演算子$\hat a_i$、$\hat a_i^{\dagger}$の標準的な交換関係$[\hat a_i,\hat a_i^{\dagger}]=1$を用いると、(5.10)を
\chfiveeq{13}
と書き換えられる。

% 原文 p.181
\hypertarget{srcpage:181}{}\label{srcpage:181}
最後の段階では［(5.12)参照］
\chfiveeq{14}
を用いた。よって、環境自己相関関数(5.6)は
\chfiveeq{15}
と書ける。ここで関数
\chfiveeq{16_17}
は、文献では一般に、それぞれ雑音核、散逸核と呼ばれる。(5.16)の第1行の波括弧$\{\cdot,\cdot\}$は、反交換子$\{\hat A,\hat B\}\equiv\hat A\hat B+\hat B\hat A$を表す。(5.16)、(5.17)の最後の行で導入した関数$J(\omega)$は
\chfiveeq{18}
と定義され、環境のスペクトル密度と呼ばれる。スペクトル密度は環境の物理的性質をまとめたものであり、デコヒーレンスの理論的・実験的研究できわめて重要な役割を果たす。環境の模型化では、離散的振動数$\omega_i$、質量$m_i$をもつ個々の振動子による記述を、環境の振動数$\omega$の連続スペクトルに対応するスペクトル密度$J(\omega)$に置き換える、環境の連続極限へ移ることが多い。

% 原文 p.182
\hypertarget{srcpage:182}{}\label{srcpage:182}
そのような連続スペクトル密度の例は、後で論じる。当面は、$J(\omega)$の元の表式(5.18)を用いる。したがって、この段階で定義(5.18)の利点は、雑音核と散逸核(5.16)、(5.17)の記法が、いくらか簡潔になるだけである。

環境自己相関関数$C(\tau)$［(5.15)参照］を決定できたので、目標とする量子Brown運動のBorn--Markovマスター方程式の導出の主要な段階は終わった。残りの導出は容易である。演算子$\hat B$、$\hat C$［(4.9a)、(4.9b)参照］は、直ちに
\chfiveeq{19a_19b}
と書ける。ここで
\begin{equation}\hat X(\tau)=e^{\mathrm i\hat H_{\mathrm S}\tau}\hat Xe^{-\mathrm i\hat H_{\mathrm S}\tau}\tag{5.20}\end{equation}
は、相互作用描像における系Sの位置演算子である。$C(\tau)$を(5.15)で与え、(5.19a)、(5.19b)を、Born--Markovマスター方程式の一般式(4.8)へ代入すると、
\chfiveeq{21}
を得る。雑音核と散逸核(5.16)、(5.17)による分解$C(\tau)=\nu(\tau)-\mathrm i\eta(\tau)$［(5.15)参照］を用いて項を整理すると、マスター方程式の最終形
\chfiveeq{22}
を得る。環境によるデコヒーレンス（および散逸）の影響は、雑音核と散逸核(5.16)、(5.17)に、系のSchrödinger描像と相互作用描像の位置演算子の積を掛けたものの積分によって、完全に記述されることがわかる。

\subsubsectionlink{中心系としての調和振動子}
マスター方程式(5.22)をさらに簡単にし、雑音項と散逸項の物理的意味を理解するため、ここからは、標準的な自己Hamiltonian
\chfiveeq{23}
をもつ調和振動子で系を表す場合に限定しよう。

% 原文 p.183
\hypertarget{srcpage:183}{}\label{srcpage:183}
相互作用描像の位置演算子$\hat X(\tau)$［(5.20)参照］の具体的な形は、系の演算子$\hat X$、$\hat P$についてのHeisenberg運動方程式を解けば容易に求まり、
\chfiveeq{24}
となる。この$\hat X(\tau)$の表式を(5.22)に代入すると、調和振動子の熱環境に線形かつ弱く結合した調和振動子の、縮約された運動を記述するBorn--Markovマスター方程式
\chfiveeq{25}
を得る。ここでは、係数$\widetilde\Omega^2$、$\gamma$、$D$、$f$を
\chfiveeq{26a_26b_26c_26d}
と定義した。これらの係数には、マスター方程式(5.25)における役割から、それぞれ特定の物理解釈がある。各項を一つずつ論じる前に、文献でよく用いられる、他の二つの同値な形へ(5.25)を書き換えよう。

第一に、(5.25)を位置表示で表すことが、しばしば便利である。基本的関係
\chfiveeq{27}
を用いると、位置基底における縮約密度行列$\rho_{\mathrm S}(X,X',t)$の運動方程式
\chfiveeq{28}
を得る。

% 原文 p.184
\hypertarget{srcpage:184}{}\label{srcpage:184}
また、位置空間の縮約密度行列$\rho_{\mathrm S}(X,X')$のWigner表示$W(X,P)$［(2.132)と2.15.2節参照］についてのマスター方程式として、(5.25)を表すこともできる。ここでは導出を示さず、結果だけを述べる\cite{ref219}。
\chfiveeq{29}
統計力学に馴染みのある読者は、(5.29)と、古典統計物理のFokker--Planck方程式との形式的な類似性に気づくかもしれない。

マスター方程式(5.25)に現れる各項の解釈を論じよう。(5.25)の右辺の最初の交換子
\chfiveeq{30}
は、物理的振動数$(\Omega^2+\widetilde\Omega^2)^{1/2}$の調和振動子の、可逆でユニタリーな運動を記述する。したがって$\widetilde\Omega^2$は、中心振動子の固有振動数$\Omega$に、温度に依存しないシフトを導入する。これは155ページで述べたLambシフトの寄与の一例である。

次に、(5.25)の項
\begin{equation}-\mathrm i\gamma[\hat X,\{\hat P,\hat\rho_{\mathrm S}(t)\}]\tag{5.31}\end{equation}
は、環境との相互作用による運動量減衰、したがって散逸を記述する。これを見るため、運動方程式(5.25)で決まる、中心振動子の運動量期待値$\langle\hat P\rangle_t$の時間発展を考えよう。トレース則$\langle\hat P\rangle=\operatorname{Tr}_{\mathrm S}(\hat\rho_{\mathrm S}\hat P)$を用いると
\chfiveeq{32}
を得る。第1項$-M(\Omega^2+\widetilde\Omega^2)\langle\hat X\rangle_t$は、Lambシフトした調和振動子ポテンシャルにおける、通常のユニタリーな発展による。当面、このユニタリーで非散逸的な項を無視しよう。すると、(5.32)より、環境との相互作用によって、中心系の運動量期待値が率$2\gamma$で指数関数的に減少することが直ちにわかる。
\chfiveeq{33}
したがって、マスター方程式(5.25)の項(5.31)は、確かに$\gamma$に比例する率の運動量減衰を記述する。この率はスペクトル密度のみに依存し、環境の温度には依存しない［(5.17)、(5.26b)参照］。

% 原文 p.185
\hypertarget{srcpage:185}{}\label{srcpage:185}
次に、(5.25)の右辺の第3項
\chfiveeq{34}
を考えよう。この項は、Lindbladの二重交換子の形(4.45)であり、環境による系の位置座標$\hat X$の監視、したがって位置基底のデコヒーレンスを記述する。この役割は、位置表示で二重交換子を書き下しても容易にわかる。
\chfiveeq{35}
読者は、右辺の項が、第3章で論じた環境散乱模型で空間的局在化を記述した表式(3.55)と、同様の形をしていることに気づくかもしれない。したがって、係数$D$と$(X-X')^2$の積は、デコヒーレンス率として働く。すなわち
\chfiveeq{36}
は、距離$|X-X'|$にわたる空間的干渉が$e$分の1に抑制される時間尺度である［(3.58)も参照］。

係数$D$は、系の運動量の二乗の期待値$\langle\hat P^2\rangle_t$の時間発展からわかるように、運動量の拡散も記述する。(5.32)を得たのと同様の簡単な計算から
\chfiveeq{37}
を得る。$D$の項だけによる$\langle\hat P^2\rangle_t$の時間依存性に注目すると、$\langle\hat P^2\rangle_t\propto Dt$という発展を得る。これは、$D$が確かに運動量空間の拡散を記述することを示す。したがって、量子Brown運動の文献では、$D$はしばしば通常拡散係数と呼ばれる。ここで「通常拡散」という用語は、$\langle\hat P^2\rangle_t$の線形な時間依存性を指す。

このような通常拡散は、古典的Brown運動の特徴であることを思い出そう。そこでは、粒子の可能な位置を記述する正規分布の分散$\sigma_x^2$が、$\sigma_x^2(t)\propto t$と増大する。実際、同じ時間依存性は量子の場合にも成り立つ。長時間極限$\gamma t\gg1$において、マスター方程式(5.25)が位置空間の分散に対して
\chfiveeq{38}
という時間依存性を与えることは容易に示せる。すなわち、位置空間アンサンブルの幅$\Delta X(t)$は、古典的Brown運動と同様、漸近的に$\Delta X(t)\propto\sqrt t$とスケールする。

% 原文 p.186
\hypertarget{srcpage:186}{}\label{srcpage:186}
実際、量子と古典のこの類似性が、もともと「量子Brown運動」という用語を選ぶ動機となった。ちなみに、$\langle\hat P^2/2M\rangle$は系の運動エネルギーなので、(5.37)は、係数$\gamma$が散逸を定量化するという先の結果(5.33)も確認する。

拡散係数としての$D$の役割は、Wigner表示のマスター方程式(5.29)からもわかる。$D$の項だけを取り出すと、発展方程式
\chfiveeq{39}
を得る。これは、一定の拡散係数$D$をもつ場合の、変数$P$についてのよく知られた拡散方程式の形である。

最後に、マスター方程式(5.25)の最終項
\chfiveeq{40}
に注目しよう。この項も、通常拡散項(5.34)と同様に、拡散とデコヒーレンスを記述する。このことは、マスター方程式のWigner形(5.29)から最もよくわかる。$f$の項だけの寄与は
\chfiveeq{41}
である。右辺に二重微分が現れる点で、この形は拡散方程式(5.39)に似ている。しかし、「通常の」拡散方程式と異なり、(5.41)の二重微分は、異なる二つの変数$X$、$P$についての1階微分で構成される。この違いから、文献では量$f$を、しばしば異常拡散係数と呼ぶ。後で示すように、多くの物理的に興味ある状況では、系の縮約された運動に対する異常拡散項(5.40)の影響は、通常拡散項(5.34)と比べて無視できるほど小さい。そのような場合、異常拡散項を除いた、マスター方程式(5.25)の簡単な形を用いてよい（5.2.4節参照）。

この模型では、通常拡散係数と異常拡散係数は、スペクトル密度$J(\omega)$と環境温度$T$の両方に依存する［(5.16)、(5.17)を用いた(5.26)参照］。これは、前述のように環境温度に依存しない、振動数シフト$\widetilde\Omega^2$と運動量減衰率$\gamma$とは対照的である。

最後に、文献では、係数(5.26)の部分的に「前Markov的」な形もしばしば論じられることに触れておく（例えば文献\cite{ref015,ref016}）。この場合、時間積分の上限をまだ無限大とせず、時刻$t$までにとどめる（163ページの議論も参照）。言い換えれば、陽な時間依存性をもつ係数$\widetilde\Omega^2(t)$、$\gamma(t)$、$D(t)$、$f(t)$、
\chfiveeq{42a_42b_42c_42d}
を考える。

% 原文 p.187
\hypertarget{srcpage:187}{}\label{srcpage:187}
「Markov的な係数」(5.26)の代わりに、これらの時間依存係数で評価したマスター方程式(5.25)は、4.2.2節で述べたRedfield型\cite{ref174,ref175}となる。すなわち、このマスター方程式は、(i)Born近似、(ii)Bornマスター方程式(4.32)の右辺の被積分関数で、遅延時刻の密度演算子$\hat\rho^{(\mathrm I)}_{\mathrm S}(t')$を$\hat\rho^{(\mathrm I)}_{\mathrm S}(t)$に置き換える、最初のMarkov近似に基づく。ただし、(4.32)の積分範囲は無限大に延ばさない。

得られたマスター方程式、すなわち係数を(5.42)で与えた(5.25)は、例えば低温環境の場合など、Markov近似が完全には適切でない状況で、量子Brown運動をより現実的に扱える\cite{ref109,ref220}。実際、係数(5.42)は、一方のMarkov的な設定［時間に依存しない係数(5.26)で評価したマスター方程式(5.25)］と、他方の量子Brown運動の厳密なマスター方程式（5.2.7節参照）の中間に位置する。後者は、ここで考える弱結合極限だけでなく、任意の系・環境相互作用の強さで有効である。

係数(5.42)の時間依存性は、文献で詳しく研究されている（例えばHu、Paz、Zhang\cite{ref219}およびPaz、Zurek\cite{ref015}の議論）。例えば、環境がOhmicなスペクトル密度をもつ場合（次の5.2.3節参照）、振動数シフト$\widetilde\Omega^2(t)$と運動量減衰係数$\gamma(t)$は、ともにゼロから始まり、環境に顕著な割合で存在する最高振動数の大きさによって定まる時間尺度で、最終値へ漸近することが示せる。拡散係数$D(t)$、$f(t)$は、温度に依存しない初期過渡応答を示した後、温度に依存する漸近値へかなり速く落ち着く。

Markov的な係数(5.26)は、(5.42)の積分で極限$t\longrightarrow\infty$を取ることに対応するため、$\widetilde\Omega^2(t)$、$\gamma(t)$、$D(t)$、$f(t)$は、これらのMarkov的係数の値へ漸近する。以下では、Markov的な場合、すなわちマスター方程式(5.25)と、係数の時間に依存しない表式(5.26)に注目する。
